\documentclass[11pt,a4paper]{article}
% -- Fonts (requires lualatex or xelatex) --
\usepackage[sfdefault,lining]{FiraSans}
\usepackage[fakebold]{firamath-otf}
\renewcommand*\oldstylenums[1]{{\firaoldstyle #1}}

% -- Layout & typography --
\usepackage{microtype}
\usepackage[margin=2.5cm]{geometry}

% -- Math --
% amssymb is omitted: firamath-otf loads unicode-math which provides
% all AMS symbols and conflicts with amssymb's definitions.
\usepackage{amsmath}

% -- Tables & figures --
\usepackage{booktabs}
\usepackage{siunitx}

% -- Symbols --
\usepackage{pifont}
\newcommand{\cmark}{\ding{51}}

% -- Colors & links --
\usepackage{xcolor}
\usepackage{hyperref}
\hypersetup{
    colorlinks=true,
    linkcolor=blue!60!black,
    urlcolor=blue!60!black,
}


\title{A Lightweight, Self-Calibrating ETA Estimator\\
       for a Bike Computer}
\author{Godot Cycling ETA -- Method D}
\date{}

\begin{document}
\maketitle

\begin{abstract}
\noindent
We describe a method for estimating remaining ride time on a GPS-capable
bike computer that requires only the route and the rider's total
weight as device-level inputs. Aerodynamics, rolling resistance, and
wind are treated as fixed assumptions, and the rider's flat-ground
cruising speed is measured from the first few minutes of the ride
itself. The physical model is a constant-power force balance that we
solve in closed form per gradient bin to produce a speed-ratio curve.
A short behavioral adjustment layer makes the curve match what real
riders actually do on climbs and descents. The route is decimated by
Visvalingam--Whyatt simplification into constant-gradient segments;
time-to-go is then a dot product between a per-gradient speed table
and a histogram of remaining route length at each gradient. A pair of
cumulative integral corrections, one for climbs and one for descents,
absorbs any steady-state bias that the physics model cannot predict
a priori. The result is an estimator with a small fixed set of inputs,
a tiny online state, and a self-calibration mechanism that converges
within the first few minutes of any ride.
\end{abstract}

\section{Problem statement}

On a device like a Karoo Hammerhead the inputs we can rely on are:
the route (a sequence of GPS coordinates with elevation), the rider's
total weight, and the 1\,\si{\hertz} stream of GPS samples arriving
as the ride unfolds. Everything else~--- power meter data, heart
rate, cadence, temperature, wind conditions, fitness history~--- may
or may not be present, and we want the estimator to be useful even
when it is not.

The ride-time prediction problem is then: given the route geometry
ahead, the weight and (implicitly) recent pace of the rider, predict
the remaining time to the finish as accurately as possible at every
1\,\si{\hertz} tick. A useful prediction is one whose error is
\emph{small and stable}: the rider wants a number they can trust in
planning decisions (when to eat, when to push, when to turn around)
rather than one that swings wildly from minute to minute.

We describe a method that achieves this with a physics model of
modest complexity, a compact route representation, and a single
observation-based correction mechanism. We call it Method~D within
the broader Godot project, where it was identified as the best
clean-architecture alternative to the static-prior baseline across a
corpus of 290 historical rides.


\section{Why the moving average is hard to beat}
\label{sec:baseline}

The simplest ETA estimator is the expanding moving average: divide
total distance covered by total moving time elapsed, and project the
remainder at that pace. It requires no model, no route knowledge, and
no parameters. Given a sufficiently long ride with roughly steady
effort and no major state changes, it converges to the correct ETA by
construction --- once enough of the ride has been observed, the
expanding average \emph{is} the rider's true pace, and the fraction
of route remaining shrinks to zero.

Beating this baseline is surprisingly difficult. Any model-based
estimator risks being \emph{wrong in a way the moving average is
not}: it introduces assumptions (physics, rider behavior,
environmental constants) that, if misspecified, produce systematic
bias the moving average would never have. The only scenario where a
model can reliably outperform is when the future route differs from
the past route in a way the model can predict but the moving average
cannot.

Gradient is exactly such a difference. A rider who has averaged
\SI{28}{\km\per\hour} over 50\,\si{\km} of flat road will not
sustain that pace on the \SI{10}{\km} climb ahead --- but the moving
average does not know a climb is coming. A gradient-aware model does,
and this is the single largest source of improvement over the baseline
in our experiments: on mountain rides the moving average produces a
MAE of 29.9\,\si{\minute}, while the physics model with split
corrections achieves 14.3\,\si{\minute} (Table~\ref{tab:buildup}).

Other state changes --- bonking, dropping out of a fast group, a
sudden headwind, mechanical problems --- are equally disruptive to
ETA but fundamentally unpredictable. We do not model them explicitly.
Instead, the cumulative integral corrections
(Section~\ref{sec:integral}) absorb whatever steady bias these events
introduce, converging over the same timescale the moving average
would. The corrections are not faster than the moving average at
adapting to surprises; they are better at \emph{not needing to adapt}
to the one effect we can predict: gradient.

We can make this concrete. Suppose a rider has averaged speed
$v_{\text{avg}}$ over the first fraction $f$ of a ride of total
distance~$D$. The moving average predicts the remaining
$(1 - f) \cdot D$ at $v_{\text{avg}}$. If the remaining portion's
true average speed is $v_{\text{true}} = \rho \cdot v_{\text{avg}}$
(where $\rho < 1$ means the rider is slower, e.g.\ climbing), the
absolute ETA error is
\begin{equation}
\label{eq:ma-error}
\varepsilon \;=\; (1 - f) \cdot \frac{D}{v_{\text{avg}}}
  \cdot \Big(1 - \frac{1}{\rho}\Big).
\end{equation}
Two features are worth noting. First, the \emph{relative} error
$\varepsilon / T_{\text{remaining}} = \rho - 1$ does not depend on
$f$~--- the moving average is always wrong by the same fraction of
the remaining time. Second, the \emph{absolute} error shrinks with
$(1 - f)$ because there is less distance left to be wrong about.
This is why the moving average converges: not because it gets smarter
about the terrain ahead, but because the window over which it is
wrong shrinks to zero.

Table~\ref{tab:ma-error} shows the absolute error for a
\SI{5}{\hour} ride at $v_{\text{avg}} = \SI{29}{\km\per\hour}$, for
several speed ratios $\rho$ and ride fractions $f$. A rider who has
averaged \SI{29}{\km\per\hour} over 50\,\% of a ride and faces a
remaining portion at $\rho = 0.5$ (a steep sustained climb, roughly
\SI{14.5}{\km\per\hour}) will see a moving-average ETA that is
\SI{2.5}{\hour} too short. Even at $\rho = 0.7$ (a moderate climb)
and 70\,\% complete, the error is still 37\,\si{\minute}.

\begin{table}[ht]
\centering
\caption{Absolute ETA error of the moving average (minutes) on a
\SI{5}{\hour} ride when the remaining portion differs from the
historical pace by a factor~$\rho$. Negative values indicate
underestimated remaining time (arrival later than predicted).}
\label{tab:ma-error}
\sisetup{round-mode=places, round-precision=0}
\begin{tabular}{@{} l S[table-format=-3.0] S[table-format=-3.0] S[table-format=-3.0] S[table-format=-3.0] S[table-format=-3.0] @{}}
\toprule
& \multicolumn{5}{c}{Ride fraction completed ($f$)} \\
\cmidrule(lr){2-6}
{$\rho = v_{\text{true}}/v_{\text{avg}}$}
  & {10\,\%} & {30\,\%} & {50\,\%} & {70\,\%} & {90\,\%} \\
\midrule
0.5 (steep climb)     & -270 & -210 & -150 & -90  & -30 \\
0.6                   & -180 & -140 & -100 & -60  & -20 \\
0.7 (moderate climb)  & -116 & -90  & -64  & -39  & -13 \\
0.8                   & -68  & -53  & -38  & -23  & -8  \\
0.9 (rolling)         & -30  & -23  & -17  & -10  & -3  \\
1.2 (tailwind/descent)& 45   & 35   & 25   & 15   & 5   \\
\bottomrule
\end{tabular}
\end{table}

A gradient-aware estimator sidesteps this entirely. Because it knows
the remaining route contains a climb at gradient~$g$ and has a ratio
table $r(g)$ that predicts the speed reduction, it does not need to
wait for the rider to actually slow down before adjusting the
prediction. The error in equation~\eqref{eq:ma-error} is replaced by
the much smaller residual between the physics model's ratio and the
rider's true ratio at that gradient --- typically a few percent rather
than a factor of two.

This asymmetry shapes the architecture. Each layer we add ---
gradient ratios, behavioral adjustments, relevance filtering, an
empirical power curve --- is designed to \emph{augment} the
baseline's natural convergence, not replace it. When a layer has no
information to contribute (no gradient on a flat ride, no power meter
on the handlebars), it falls back to identity: the correction stays
at~1, the ratio table collapses to $r(g) = 1$, and the estimator
degenerates to the moving average. The goal is to never be worse than
the baseline, and to be substantially better when the route gives us
something to work with.


\section{Handling pauses}
\label{sec:pauses}

A rider who stops for coffee, waits at a traffic light, or fixes a
flat tire is not generating useful speed observations. How the
estimator treats these pauses has a large effect on prediction
stability.

The naive approach is to use \emph{total elapsed time} as the
denominator: $v_{\text{avg}} = d / t_{\text{elapsed}}$. This makes
the average speed drop during every pause, and the longer the pause,
the lower the average, and the higher the predicted remaining time.
The behavior mirrors the gradient problem in
Table~\ref{tab:ma-error}: the rider's speed appears to drop (to
zero, in fact), and the estimator responds by inflating the ETA
for the entire remaining route. A 20-minute coffee stop on a
3-hour ride can add 30+ minutes to the prediction, all of which
vanishes the moment the rider clips back in and the average
recovers. The result is an ETA that swings wildly around every stop.

The better approach, and the one we use throughout, is to track
\emph{moving time} only. Rows where the rider is stationary (speed
below a threshold, or a gap in timestamps exceeding a few seconds)
are flagged as paused and excluded from all accumulators: the
expanding average, the integral corrections, and the histogram
update. The estimator predicts remaining \emph{moving time}, which
is unaffected by how long the rider has been standing still.

To convert to a wall-clock ETA, we simply add the observed pause
time so far to the moving-time prediction:
\begin{equation}
\text{ETA}_{\text{wall}} \;=\;
  \text{ETA}_{\text{moving}} \;+\; t_{\text{paused}}.
\end{equation}
This has two benefits. First, the remaining moving time is stable
across pauses --- a coffee stop does not cause the prediction to
spike. Second, the split between moving time and pause time is
transparent to the rider: they know better than any model whether
another stop is coming (a planned lunch) and know as little as the
model about unplanned stops (a mechanical). By separating the two,
the estimator handles what it can predict (moving pace) and leaves
what it cannot (future pauses) to the rider's own judgment.


\section{Physics model: from power to speed}

\subsection{Force balance on a slope}

For a cyclist travelling at steady-state speed $v$ on a road of slope
angle $\theta$, balancing power output against resistive forces gives
\begin{equation}
\label{eq:force-balance}
P(\theta) \;=\; \tfrac{1}{2}\,\rho\,C_d A\,(v + v_w)^2\,v
              \;+\; C_{rr}\,m\,g\,v\cos\theta
              \;+\; m\,g\,v\sin\theta,
\end{equation}
where $P$ is the rider's power output at that gradient, $\rho$ is air
density, $C_dA$ is the aerodynamic drag-area coefficient, $C_{rr}$ is
the rolling resistance coefficient, $m$ is the total rider-plus-bike
mass, $g$ is gravitational acceleration, and $v_w$ is the effective
headwind component along the direction of travel. The three terms on
the right are aerodynamic drag, rolling friction, and gravitational
work against the slope, respectively.

On flat ground ($\theta = 0$) we can invert
\eqref{eq:force-balance} to get the rider's \emph{flat-ground power
budget}
\begin{equation}
\label{eq:p-flat}
P_{\text{flat}} \;=\; \tfrac{1}{2}\,\rho\,C_d A\,(v_{\text{flat}} + v_w)^2\,v_{\text{flat}}
                 \;+\; C_{rr}\,m\,g\,v_{\text{flat}},
\end{equation}
given a reference flat-ground cruising speed $v_{\text{flat}}$.
Equation~\eqref{eq:p-flat} is the link between a quantity the device
can \emph{observe} ($v_{\text{flat}}$) and the latent quantity the
physics model needs ($P_{\text{flat}}$).

\subsection{Solving for $v$ at each gradient}

Taken as a function of $v$, equation~\eqref{eq:force-balance} is a
cubic, so for a given power budget $P$ and gradient $\theta$ it has a
closed-form real root via Cardano's formula. The root gives the speed
the rider would sustain at that gradient under the assumption that
their power output equals $P$. Iterating this over a grid of integer
percent gradients $g \in \{g_{\min}, \dots, g_{\max}\}$ produces a
\emph{speed-ratio table}
\begin{equation}
r(g) \;=\; \frac{v(g)}{v_{\text{flat}}},
\qquad r(0) = 1,
\end{equation}
whose shape encodes how much slower (or faster) the rider moves at
each gradient relative to their flat-ground pace. This is the only
route-independent knowledge the estimator needs about the rider.

\subsection{Behavioral adjustments}

A constant-power assumption is physically honest but biologically
wrong: real riders do not put out identical wattage on a 10\,\% climb
and on a flat road. They push harder on climbs up to a physiological
ceiling and ease off on descents where gravity is doing most of the
work. We introduce two small lumped parameters to capture this.

Let $P_{\max}$ be the rider's ceiling (taken as a small multiple of
$P_{\text{flat}}$, or derived from an FTP estimate if one is
available). On climb gradients we replace $P$ in
\eqref{eq:force-balance} with a gradient-scaled power budget
\begin{equation}
\label{eq:climb-power}
P_{\text{climb}}(g) \;=\; \min\!\big(P_{\max},\;
  P_{\text{flat}} \cdot (1 + \alpha \cdot g)\big),
\qquad g > 0,
\end{equation}
where $\alpha$ is the \emph{climb effort} parameter that controls how
aggressively the rider ramps into their ceiling. On descents we apply
an exponential decay
\begin{equation}
\label{eq:descent-power}
P_{\text{descent}}(g) \;=\; P_{\text{flat}} \cdot \exp(-\kappa \cdot |g|),
\qquad g < 0,
\end{equation}
modelling the fact that riders progressively reduce effort as the
slope does more of the work and cornering and confidence become the
limiting factors.

With these substitutions, the cubic in \eqref{eq:force-balance} is
solved once per gradient bin, producing a ratio table $r(g)$ whose
shape reflects both physics and rider behavior. Empirically, this
shape is much more stable across operating points than the pure
constant-power version: a 2$\times$ change in $v_{\text{flat}}$
produces roughly 1.2--1.3$\times$ shape distortion in the behavioral
table versus 2.4--2.7$\times$ in the constant-power table. The
behavioral model is therefore less sensitive to $v_{\text{flat}}$
misspecification during the warmup phase described below.

\subsection{Route-adaptive behavioral parameters}
\label{sec:route-adaptive}

The behavioral parameters $\alpha$ and $\kappa$ encode a pacing
assumption: how hard does a typical rider push on climbs, and how much
do they back off on descents? But pacing is not a fixed trait --- it
depends on what the rider knows about the route ahead. A rider facing
a single 2\,\si{\km} hill on a flat century will stand up and hammer
it; the same rider facing a 30\,\si{\km} Alpine pass will settle into
a sustainable rhythm well below their short-effort ceiling. Descents
follow a similar pattern: mountain descents are steeper, longer, and
more technical, leading to more braking than the gentle slopes of a
rolling ride.

We can exploit this by selecting $\alpha$ and $\kappa$ based on the
route's difficulty classification, which is known at ride start from
the VW segments. The integral corrections' final values across our
corpus reveal the systematic bias at each difficulty tier: the
parametric model overpredicts climb speed by $\sim$20\,\% on
mountains ($\Delta_c \approx 0.80$), $\sim$7\,\% on hills
($\Delta_c \approx 0.93$), and $\sim$0\,\% on minor hills
($\Delta_c \approx 1.00$). Descent speed is overpredicted by
6--15\,\% across all categories ($\Delta_d \approx 0.85$--$0.94$).

Reducing $\alpha$ for harder routes and reducing $\kappa$ globally
brings the physics model closer to reality \emph{before the
corrections have any data}. The practical effect is largest on
front-loaded routes, where the corrections have not yet converged and
the raw physics model carries the full prediction burden. On our
corpus, route-adaptive parameters reduce mountain MAE from
11.6 to 9.8\,\si{\minute} --- matching the empirical power curve's
result without requiring a power meter or rider history.

\subsection{Cross-ride parameter learning}
\label{sec:cross-ride}

The route-adaptive defaults above are generic: they reflect average
rider behavior at each difficulty tier. A better approach is to let
each rider's own correction history tune the parameters
automatically.

After each ride the final correction values $\Delta_c$ and $\Delta_d$
are a direct measurement of how wrong the physics model was. A
correction of $\Delta_c = 0.92$ means the model overpredicted climb
speed by 8\,\%. This information can be fed back into the behavioral
parameters via a simple proportional update:
\begin{align}
\label{eq:alpha-update}
\alpha_{\text{new}} &\;=\; (1 - \lambda)\,\alpha_{\text{old}}
  \;+\; \lambda\,\alpha_{\text{ride}} \cdot \Delta_c, \\
\kappa_{\text{new}} &\;=\; (1 - \lambda)\,\kappa_{\text{old}}
  \;+\; \lambda\,\kappa_{\text{ride}} / \Delta_d,
\end{align}
where $\lambda \approx 0.2$ gives a 5-ride half-life and
$\alpha_{\text{ride}}$, $\kappa_{\text{ride}}$ are the values used on
the just-completed ride. If $\Delta_c < 1$ (rider climbed slower than
predicted), $\alpha$ decreases; if $\Delta_d < 1$ (rider descended
slower than predicted), $\kappa$ increases. Over 10--15 rides the
parameters converge to rider-specific values and the physics model
starts each ride with corrections near~1.

The entire rider state is two floating-point numbers --- $\alpha$ and
$\kappa$ --- which can be persisted across rides with negligible
storage. A device-side implementation does not need access to
historical ride files: it only needs to read the final $\Delta_c$ and
$\Delta_d$ from the ride that just ended and update two scalars. This
makes the learning mechanism compatible with bike computers where
extension software cannot access past ride data but can persist its
own state from ride to ride.

The update is conservative by design. A single outlier ride (a bonk,
a race effort, a mechanical) nudges the parameters by at most
$\lambda \approx 20\,\%$ of the correction, and the EWMA smooths it
out over subsequent normal rides. No explicit outlier detection is
needed: the same mechanism that makes the integrals robust
within a ride makes the parameter update robust across rides.

\subsection{What we assume}

The physics model requires the following inputs:
\begin{itemize}
  \item The total rider-plus-bike mass $m$ (measured).
  \item The flat-ground cruising speed $v_{\text{flat}}$ (measured from
        the ride; see Section~\ref{sec:warmup}).
  \item Fixed assumed values for $\rho$, $C_dA$, $C_{rr}$,
        $\alpha$, $\kappa$, $P_{\max}/P_{\text{flat}}$ and effective
        headwind $v_w$.
\end{itemize}

The fixed assumptions deserve a brief comment. We do \emph{not} claim
these values are accurate for any particular rider: a given rider's
$C_dA$ may be off by $\pm 20\,\%$ from the default; the rolling
resistance of their tires may differ; the wind on any given day is
entirely unmodeled. We only claim that these values are approximately
correct, that deviations from them are slowly varying and similar
across gradient bins, and that the integral correction described in
Section~\ref{sec:integral} will absorb the scale of whatever the
fixed assumptions got wrong. The physics model does not need to be
correct in absolute terms. It only needs to produce a
\emph{shape-correct} ratio table~--- one where the \emph{relative}
difficulty of a 10\,\% climb versus a 5\,\% climb is approximately
right~--- because shape errors are what the integral correction is
\emph{unable} to fix.

\section{Route representation}
\label{sec:route}

\subsection{Visvalingam--Whyatt decimation}

A raw GPS route is a noisy polyline with many thousands of vertices,
each contributing a slightly different gradient. For ETA purposes we
do not care about the fine detail: we want a compact set of
constant-gradient segments that captures the essential character of
the route. We use the Visvalingam--Whyatt (VW) line simplification
algorithm, which iteratively removes the vertex whose triangle of
elimination has the smallest area, until a minimum-area threshold
(typically a few hundred square meters of horizontal$\times$vertical
tolerance) is reached. The result is a polyline with $S$ vertices and
$S - 1$ segments, each having a constant gradient.

For a typical 5\,\si{\hour} rolling ride, VW produces on the order of
$S \approx 100$--$200$ segments; a flat commute gives fewer; an Alpine
climb more. The set of segments is the only route-dependent data the
estimator needs at runtime.

\subsection{From segments to a gradient histogram}

Given the segment list $\{(L_j, g_j)\}_{j=1}^{S}$, where $L_j$ is the
segment length in meters and $g_j$ is its gradient in percent,
we form a histogram over gradient bins
\begin{equation}
\label{eq:hist}
H(g) \;=\; \sum_{j\,:\,\text{bin}(g_j) = g} L_j,
\end{equation}
which records the total route length at each gradient bin. A 50\,\si{\km}
ride might have $H(0) = \SI{32}{\km}$, $H(3) = \SI{5}{\km}$,
$H(5) = \SI{3}{\km}$, and so on for the gradients that appear in the
route. The histogram has at most $|g_{\max} - g_{\min}| + 1 \approx 40$
entries and compresses the route geometry into a tiny fixed-size
data structure.

\section{Time-to-go as a dot product}
\label{sec:dotproduct}

The remaining time from any point on the route to the end, assuming
the rider is moving at $v_{\text{flat}} \cdot r(g)$ on every segment,
is
\begin{equation}
\label{eq:ttg}
\text{TTG} \;=\; \sum_{j \text{ ahead}} \frac{L_j}{v_{\text{flat}} \cdot r(g_j)}.
\end{equation}
Regrouping this sum by gradient bin,
\begin{equation}
\label{eq:ttg-binned}
\text{TTG} \;=\; \frac{1}{v_{\text{flat}}} \sum_{g} \frac{H^{\ast}(g)}{r(g)}
\;=\; \frac{1}{v_{\text{flat}}} \,\mathbf{H}^{\ast} \cdot \mathbf{r}^{-1},
\end{equation}
where $H^{\ast}(g)$ is the histogram over only the part of the route
that is still ahead of the rider, and $\mathbf{r}^{-1}$ is the
vector $(1/r(g))_g$. Equation~\eqref{eq:ttg-binned} is a dot product
between two small vectors of fixed length.

The representation is incremental. As the rider advances by a
distance $\Delta d$ in the current 1\,\si{\hertz} tick,
$H^{\ast}(g_{\text{current}})$ is decremented by $\Delta d$, and the
dot product is re-evaluated in constant time. When $\Delta d$ spans a
VW segment boundary (which happens rarely at typical cycling speeds
and short tick intervals) the decrement is split across the bins of
the segments involved. Route state ($\mathbf{H}^{\ast}$) and rider
state ($\mathbf{r}^{-1}, v_{\text{flat}}$) are fully decoupled:
updating one does not require touching the other.

\section{Cumulative integral correction}
\label{sec:integral}

The physics model gives a shape for $r(g)$; it does not give an
absolute pace that matches reality for any specific rider. Even with
correct $v_{\text{flat}}$, a rider may climb harder than the model
predicts (rider is better than average on hills), descend more
conservatively (rider does not like cornering at speed), or simply
run into a sustained headwind the model has no way to know about.

We handle this with two scalar correction factors $\Delta_c$ and
$\Delta_d$, one for climb rows and one for descent rows, each a ratio
of cumulative predicted time to cumulative actual time over the ride
so far:
\begin{align}
\Delta_c(t)
&\;=\;
\frac{\sum_{i \in \text{climb rows},\; \tau_i \le t} dd_i / \big(v_{\text{flat}} \cdot r(g_i)\big)}
     {\sum_{i \in \text{climb rows},\; \tau_i \le t} dt_i}, \\[1ex]
\Delta_d(t)
&\;=\;
\frac{\sum_{i \in \text{descent rows},\; \tau_i \le t} dd_i / \big(v_{\text{flat}} \cdot r(g_i)\big)}
     {\sum_{i \in \text{descent rows},\; \tau_i \le t} dt_i}.
\end{align}
A value of $\Delta_c = 1$ means the climb predictions match reality.
$\Delta_c < 1$ means the rider climbs faster than predicted (the
cumulative predicted time is less than the cumulative actual time).
$\Delta_c > 1$ means the rider climbs slower than predicted.

At query time, the ETA formula splits the TTG into its climb and
descent contributions and divides each by the corresponding integral:
\begin{equation}
\label{eq:eta}
\text{ETA}(t) \;=\;
 \frac{1}{v_{\text{flat}}} \Big[\frac{1}{\Delta_c(t)}\,\mathbf{H}^{\ast}_c \cdot \mathbf{r}_c^{-1}
                   \;+\; \frac{1}{\Delta_d(t)}\,\mathbf{H}^{\ast}_d \cdot \mathbf{r}_d^{-1}\Big],
\end{equation}
where the subscripts $c$ and $d$ denote the climb-only and
descent-only slices of the histogram and the ratio vector. The two
integrals act as scale corrections that absorb any steady bias between
what the physics model predicts and what the rider actually achieves,
independently for climbs and descents.

Algebraically, $\Delta_c \cdot v_{\text{flat}}$ is the back-derived
$v_{\text{flat}}$ inferred from climb rows alone: the rider's actual
pace on climbs, expressed in flat-equivalent units via the ratio
table. The same for $\Delta_d$ and descent rows. The ETA formula can
be read as ``weight the remaining route length by its gradient shape,
then scale by the rider's observed pace on that kind of terrain so
far.'' The $v_{\text{flat}}$ prior drops out algebraically whenever
both the histogram and the integral scale with it in the same way,
which is one reason the method is robust to a mildly wrong flat-ground
prior.

\section{Warmup and one-shot recalibration}
\label{sec:warmup}

The integral correction above is well-defined once the ride has
produced at least a few climb and descent observations. During the
first few minutes, though, the ratio table is built from a
\emph{guess} at $v_{\text{flat}}$~--- whatever the config supplies or
the last ride's value. If the guess is badly wrong, the ratio table's
shape is also off, and the integrals will spend the first half of the
ride compensating for a miscalibrated physics model rather than
correcting for genuine rider behavior.

Method~D addresses this with a single-shot recalibration step. For
the first
\begin{equation}
\tau_{\text{warmup}} \;=\;
\mathrm{clip}\!\big(\phi \cdot \hat{T}_{\text{total}},\;
  \tau_{\min},\; \tau_{\max}\big)
\end{equation}
seconds of moving time (a fixed fraction $\phi \approx 0.10$ of the
estimated total ride time $\hat{T}_{\text{total}}$, clipped into
$[\SI{3}{\minute}, \SI{10}{\minute}]$), the estimator does not yet
accumulate integrals. Instead, it runs a \emph{prior-free}
$v_{\text{flat}}$ observer in parallel: a running mean of
$v_i / r(g_i)$ over all valid moving rows, which back-derives the
rider's flat-equivalent pace by inverting the ratio table. This
observer is unbiased in expectation (it does not depend on the guess
at $v_{\text{flat}}$), but it is noisy with few samples and needs
some data to settle.

At the warmup boundary, the estimator reads off the current value
$\hat{v}_{\text{flat}}$ from the prior-free observer and \emph{rebuilds
the ratio table} at that new operating point, using
equation~\eqref{eq:force-balance} with the updated $v_{\text{flat}}$.
From then on the rebuilt table is used for all further predictions,
and the climb/descent integrals begin accumulating against the
rebuilt prediction. This is where the ``D'' comes from in the
variant's internal name: compared to the static-prior baseline, the
differentiation is \emph{one shot of calibration up front}.

After the warmup, $v_{\text{flat}}$ is frozen and the integrals do
the rest of the work. The frozen value is not revisited. This is a
deliberate choice: the two integrals have infinite memory (they are
cumulative, not windowed), and letting $v_{\text{flat}}$ drift on top
of them would introduce a second feedback loop whose dynamics would
interact in hard-to-predict ways. Freezing $v_{\text{flat}}$ post-
warmup gives the integrals a stable target to converge to, and the
integrals handle all within-ride variation from there.

Empirically, this single recalibration is enough to make the
estimator robust to a $\pm$30\,\% error in the initial
$v_{\text{flat}}$ guess: when the baseline's MAE swings by
$\sim$1.2\,\si{\minute} as we move the prior from \SI{20}{\km\per\hour}
to \SI{40}{\km\per\hour}, Method~D's MAE swings by
$\sim$0.2\,\si{\minute}.

\section{Relevant-segment filtering}
\label{sec:relevant}

The integral corrections $\Delta_c$ and $\Delta_d$ absorb
\emph{average} bias across all climb and descent observations,
respectively. On a ride that mixes short punchy walls with long
sustained climbs, the integral learns a weighted average of two
different regimes: a rider who hammers a 200\,\si{\metre} muurke at
double their flat-ground power, and the same rider pacing a
10\,\si{\km} climb at 1.3$\times$. Because the integral is
cumulative, the muurke's contribution is proportional to its
duration---typically 1--2\,\si{\minute}---and is eventually washed
out by longer observations. But if the muurke appears early in the
ride and the first long climb does not come for another hour, the
integral carries a biased $\Delta_c$ across that entire gap.

We address this by \emph{pre-classifying} VW segments at ride start
into relevant and non-relevant. A segment is \emph{relevant} when its
predicted traversal time at the prior $v_\text{flat}$ exceeds a
threshold $\tau_\text{min}$ (default 120\,\si{\second}):
\begin{equation}
\label{eq:relevant}
\text{relevant}(j)
\;\iff\;
\frac{L_j}{v_\text{flat} \cdot r(g_j)}
\;>\; \tau_\text{min}.
\end{equation}
This classification uses only the route geometry and the prior
$v_\text{flat}$; it is computed once and fixed for the ride.

The integral corrections then accumulate only over rows whose
containing VW segment is relevant:
\begin{equation}
\Delta_c(t) \;=\;
\frac{\displaystyle\sum_{\substack{i \in \text{climb},\;\tau_i \le t \\ \text{relevant}(j_i)}}
  dd_i / \bigl(v_\text{flat} \cdot r(g_i)\bigr)}
     {\displaystyle\sum_{\substack{i \in \text{climb},\;\tau_i \le t \\ \text{relevant}(j_i)}} dt_i},
\end{equation}
and analogously for $\Delta_d$. Non-relevant segments are still
predicted---they contribute to the TTG via
equation~\eqref{eq:ttg}---but their samples do not feed back into
the correction state. On a ride with no relevant segments (a typical
flat Belgian day) the integrals simply stay at~1 and the estimator
degenerates to the uncorrected physics model, which is the right
behavior when there is nothing material to learn from.

On a corpus of 290~rides partitioned by PCS-inspired climb
difficulty, this filter reduces MAE on mountain rides by
$\sim$19\,\% (from 14.3 to 11.6\,\si{\minute}) while leaving flat
rides unchanged. The threshold $\tau_\text{min} = 120$\,\si{\second}
was selected by a sweep over $\{60, 120, 180, 300, 600\}$; it
corresponds to the approximate boundary between a punchy wall the
rider can sustain above threshold and a paced effort where riding
behavior dominates sprint dynamics.


\section{Power-informed behavioral model}
\label{sec:power}

\subsection{Replacing the parametric model with measurement}

The behavioral adjustments in equations~\eqref{eq:climb-power} and
\eqref{eq:descent-power} introduce two free parameters ($\alpha$ and
$\kappa$) that control how power varies with gradient. Their values
are hand-tuned to a generic rider. When a power meter is available,
we can replace the parametric family entirely with a measured
\emph{power curve}
\begin{equation}
\label{eq:pi}
\pi(g) \;=\; \frac{P(g)}{P_\text{flat}},
\qquad \pi(0) = 1,
\end{equation}
computed as the median power ratio per integer gradient bin, either
from the rider's historical rides or from the current ride as
observations accumulate.

The ratio table $r(g)$ is then produced by the same cubic solver
as before, but with the behavioral power budget at each gradient
replaced by the empirical one:
\begin{equation}
\label{eq:r-empirical}
r(g) \;=\; \frac{v(g)}{v_\text{flat}},
\qquad\text{where } v(g) \text{ solves~\eqref{eq:force-balance} at }
P = P_\text{flat} \cdot \pi(g).
\end{equation}
The physics (drag, rolling resistance, gravity, headwind) is
unchanged; only the \emph{rider intent} input to the cubic is swapped
from a parametric assumption to a lookup. Gradient bins for which
$\pi(g)$ has not been observed fall back to the parametric behavioral
model, so the empirical curve strictly refines the physics rather
than replacing it.

\subsection{Decomposition: physics residual and behavioral signal}

With measured power on every sample, the observed speed ratio
decomposes as
\begin{equation}
\label{eq:decomposition}
\underbrace{\frac{v_\text{obs}}{v_\text{flat}}}_{\text{observed ratio}}
\;=\;
\underbrace{\frac{v_\text{obs}}{v_\text{phys}(P_\text{obs}, g)}}_{\text{physics residual}}
\;\cdot\;
\underbrace{\frac{v_\text{phys}(P_\text{obs}, g)}{v_\text{flat}}}_{\text{ratio at measured power}},
\end{equation}
where $v_\text{phys}(P, g)$ is the steady-state speed from the cubic
at the \emph{measured} power $P_\text{obs}$. The first factor is a
pure physics residual (CdA, $C_{rr}$, wind errors); the second
encodes rider intent via the cubic. Empirically, the physics residual
is approximately flat across climb gradients ($\sim$0.93--0.96),
confirming that rider intent---not physics error---is the dominant
source of shape mismatch in the ratio table. On descents the residual
drops sharply (to $\sim$0.66 at $-8$\,\%) reflecting braking and
cornering, which is a speed-level phenomenon that the power channel
cannot resolve.

This decomposition motivates the asymmetric architecture: on climbs,
the empirical power curve $\pi(g)$ captures rider intent and the
integral correction $\Delta_c$ absorbs the small, flat physics
residual. On descents, $\pi(g)$ is less useful (power does not
predict speed well when braking dominates), so $\Delta_d$ carries
the full correction as before.

\subsection{Three tiers of availability}
\label{sec:power-tiers}

The estimator operates in three modes depending on what data is
available:
\begin{enumerate}
\item \textbf{No power meter.}
  The parametric behavioral model
  (equations~\ref{eq:climb-power}--\ref{eq:descent-power}) provides
  $r(g)$. This is the default and requires no rider history. The
  relevant-segment filter (Section~\ref{sec:relevant}) and the
  split integral corrections provide self-calibration from speed
  observations alone.
\item \textbf{Power meter, no history.}
  The empirical $\pi(g)$ is estimated from the current ride as
  gradient bins accumulate samples ($\ge$30 per bin). Bins without
  enough data fall back to the parametric model. On hilly rides this
  produces a usable curve within the first 30--60\,\si{\minute}; on
  mountain rides with sparse steep-gradient bins the fallback
  dominates and the benefit is smaller.
\item \textbf{Power meter with rider history.}
  A cached $\pi(g)$ profile, computed from the rider's historical
  rides, is loaded at ride start. All bins are populated and the
  ratio table is rider-specific from the first prediction. The
  current ride's power observations can further refine the profile
  via a trust-weighted blend.
\end{enumerate}

On a corpus of 290~rides, the cached profile reduces mountain-ride
MAE by 30\,\% versus the unfiltered baseline (from 14.3 to
9.9\,\si{\minute}) and the per-ride mode by 25\,\% (to
10.7\,\si{\minute}). On flat rides the improvement is smaller but
consistent ($\sim$3\,\% MAPE reduction), confirming that the
empirical curve is a strict improvement over the parametric model
even on terrain where gradient-awareness matters least.


\section{What we do not model}

It is worth being explicit about what the estimator deliberately
leaves out, and why.

\paragraph{Fatigue.} We do not model progressive decline in rider
power output as the ride proceeds. If the rider slows down in the
second half of the ride because they are tired, the integrals will
drift smoothly toward the new equilibrium over the following
$\sim$30\,\si{\minute} (set by the rate at which new observations
accumulate relative to old ones in the expanding-mean integral).
For rides where fatigue is mild, this drift is accurate enough. For
rides where fatigue is extreme (ultra-distance events where the rider
bonks) the integral adapts more slowly than reality and the estimator
will underpredict the remaining time during the transition.

\paragraph{Wind.} The effective headwind $v_w$ in
equation~\eqref{eq:force-balance} is a fixed assumption. If the real
wind differs from the assumption, the entire ratio table is
distorted, but predominantly uniformly: a tailwind shifts all ratios
up, a headwind shifts them down. The climb/descent integrals absorb
this as a scale error without needing to know why it happened.

\paragraph{Temperature, humidity, altitude.} These affect air density
$\rho$ slightly and the rider's power output via physiology. Both
effects are second-order relative to gradient and are treated as
constants.

\paragraph{Group dynamics.} Drafting in a peloton reduces aerodynamic
drag by a factor that depends on position in the group and the
group's speed, which is a function of minute-by-minute tactical
decisions. We do not model this. If the rider joins a fast group,
their effective pace increases in a way the climb/descent integrals
will eventually catch up with, but probably not before the benefit
evaporates.

\paragraph{Traffic and stops.} These are handled upstream of the
estimator: the GPS track is pre-processed to flag paused rows, and the
estimator's integrals and TTG computation skip them. The wall-clock
ETA is adjusted separately to reflect the time spent paused so far.

\paragraph{Mid-ride state changes.} We do not attempt to detect
abrupt shifts in rider physiology or conditions during the ride. We
experimented with several mechanisms (divergence monitoring, pause-
triggered resets, change-point detection) and found empirically that
none of them produced reliable gains on a corpus of 311 real rides.
The integrals' own drift rate is sufficient to catch most
state-change effects over the typical timescale of a ride, and
attempts to detect state changes explicitly either produced no
improvement or introduced spurious resets that hurt more than they
helped.

\section{Why the unmodeled effects are tolerable}

The estimator's central claim is that, as long as ride duration is
sufficiently long relative to the integral correction's effective
time constant, the cumulative integrals will dominate any steady
bias introduced by the fixed physics assumptions. We can make this
precise as follows.

Suppose the physics model underestimates rider pace by a factor of
$(1 + \epsilon)$, uniformly across all gradients, for reasons we do
not know and cannot correct a priori. After $n$ moving observations
have accumulated in the integrals, both $\Delta_c$ and $\Delta_d$
will have converged to values that absorb this factor, and the ETA
formula \eqref{eq:eta} will scale by the same $(1 + \epsilon)$ to
match reality. The bias is gone. The only lingering error is on the
rows \emph{before} the integrals have enough samples to converge~---
which is the warmup period we already handle with the one-shot
recalibration.

For non-uniform bias~--- say the physics model is right about climbs
but wrong about descents, or vice versa~--- the split into climb and
descent integrals gives each of the two terrain classes its own
correction. The only bias the method cannot absorb is a bias that
varies non-uniformly \emph{within} the climb class (rider climbs 10\%
harder than predicted on 5\,\% gradients but only 5\% harder on
10\,\% gradients) or within the descent class. Such biases would
correspond to genuine shape errors in the ratio table, and the
behavioral tweaks in Section~2 are what the method offers as a
mitigation. They are not perfect, but across our corpus the residual
shape error within a terrain class is small enough that the integrals
do most of the work.

The upshot is that, for rides lasting at least $\sim$30 minutes and
for rider-environment mismatches that are steady over the ride,
Method~D converges on accurate predictions whose error is dominated
by the initial warmup and bounded by the quality of the behavioral
ratio table. For the rider's own repeated rides on their usual routes,
this error is typically around $\pm$6\,\si{\minute} of MAE on a 5-hour
ride, or roughly 2\,\%.

\section{State size and online cost}

For a device-side implementation, the estimator maintains:
\begin{itemize}
  \item The gradient histogram $\mathbf{H}^{\ast}$: one float per
        bin that appears in the route. Typically 20--40 floats.
  \item The inverse-ratio vector $\mathbf{r}^{-1}$: one float per
        gradient bin. Same size.
  \item The current segment index and distance-into-segment for the
        incremental update. Two scalars.
  \item The running climb and descent integrals: four scalars
        (numerator and denominator of $\Delta_c$ and $\Delta_d$).
  \item The rider parameters: $v_{\text{flat}}$, $m$, and the fixed
        physics constants. A few more scalars.
\end{itemize}
Total state size is on the order of a few hundred bytes. Per-tick
cost is a constant-time histogram update plus two constant-time
integral updates plus a dot product of size $\sim$40.

\section{Metrics}
\label{sec:metrics}

All metrics compare predicted remaining moving time to actual
remaining moving time at each 1\,\si{\hertz} tick, excluding the
first 2\,\% of ride distance (warmup). Unless stated otherwise,
results report averages across rides in each difficulty bucket.

\begin{description}
\item[MAE (Mean Absolute Error, minutes)]
  $\frac{1}{N}\sum_i |\text{ETA}_i - \text{ATA}_i|$.
  The average prediction error across all ticks. Dominated by
  mid-ride ticks where absolute remaining time is large and the
  prediction has the most room to be wrong. A constant-speed oracle
  (true whole-ride average) does \emph{not} give $\text{MAE} = 0$:
  it gives 1.8\,\si{\minute} on pancake rides and 5.6 on hills,
  because the rider's speed varies within the ride and the remaining
  portion's average speed differs from the whole-ride average at
  every moment.
\item[MPE (Mean Percentage Error, \%)]
  $\frac{1}{N}\sum_i (\text{ETA}_i - \text{ATA}_i) / \text{ATA}_i
  \times 100$. Signed: positive means the estimator overpredicts
  remaining time (rider arrives earlier than predicted); negative
  means the rider arrives later. An unbiased estimator has
  $\text{MPE} \approx 0$.
\item[MAPE (Mean Absolute Percentage Error, \%)]
  $\frac{1}{N}\sum_i |\text{ETA}_i - \text{ATA}_i| / \text{ATA}_i
  \times 100$. Normalises by remaining time, so a 5-minute error at
  2\,hours remaining (4\,\%) is weighted the same as a 1-minute error
  at 24\,minutes remaining (4\,\%). Less dominated by mid-ride ticks
  than MAE.
\item[RMSE (Root Mean Square Error, minutes)]
  $\sqrt{\frac{1}{N}\sum_i (\text{ETA}_i - \text{ATA}_i)^2}$.
  Penalises large errors more than MAE. A useful companion metric:
  if RMSE $\gg$ MAE, the estimator has occasional large spikes
  rather than uniform moderate error.
\item[Settle time (minutes)]
  The moving time until the prediction stabilises within
  $\pm$5\,\% of its final-tick value. Measures convergence speed:
  how long the rider must wait before the ETA becomes trustworthy.
\end{description}

The distinction between MAE and the oracle floor deserves emphasis.
An estimator with $\text{MAE} = 5$\,\si{\minute} on hills is not
``5\,minutes wrong on average'' in the sense a rider might expect.
It means: across all moments of all rides, the average absolute
deviation between predicted and actual remaining time is
5\,\si{\minute}. Early in the ride (when remaining time is
$\sim$3\,\si{\hour}) the error may be 10+\,\si{\minute}; near the
end it may be $<$\,1\,\si{\minute}. The MAE averages these. For a
rider making pacing decisions, the error at a specific ride fraction
(e.g.\ 50\,\%, 75\,\%) is often more relevant than the ride-average
MAE.


\section{Results}
\label{sec:results}

Table~\ref{tab:buildup} shows the effect of each architectural layer
on a corpus of 290 historical rides, partitioned by route difficulty
using a PCS-inspired climb-score classification based on the hardest
single climb: \emph{pancake} (score $< 0.5$, no meaningful gradient),
\emph{rolling} ($0.5$--$2$, minor undulations),
\emph{minor hills} ($2$--$5$, small real hills),
\emph{hills} ($5$--$50$, proper short climbs),
and \emph{mountains} ($\ge 50$, at least one sustained climb).
Each row adds exactly one feature to the previous row; the MAE
columns show mean absolute ETA error in minutes. The physics prior
rows use the parametric behavioral model without any online
correction, illustrating the gap the corrections must close.

\begin{table}[ht]
\centering
\caption{Build-up of estimator complexity. Shaded rows are oracle
bounds that cheat with end-of-ride information (see
Section~\ref{sec:oracle}). Each subsequent row adds one feature to
the row above. MAE is mean absolute ETA error in minutes (lower is
better). \emph{Behav.}: ``--'' = constant power, ``param'' =
parametric, ``emp'' = empirical $\pi(g)$. \emph{$\Delta$}: correction
channels --- ``4-dur'' = duration-bucketed (short/medium/long climb +
descent). \emph{Rel.}: relevance filter --- ``cond'' = mountains
only.}
\label{tab:buildup}
\sisetup{round-mode=places, round-precision=1}
\small
\begin{tabular}{@{} l ccccc S[table-format=2.1] S[table-format=2.1] S[table-format=2.1] S[table-format=2.1] S[table-format=2.1] @{}}
\toprule
& \multicolumn{5}{c}{Features}
& \multicolumn{5}{c}{MAE (min)} \\
\cmidrule(lr){2-6} \cmidrule(lr){7-11}
Estimator
  & {$r(g)$} & {Behav.} & {$\Delta$} & {Rel.} & {Power}
  & {Pancake} & {Rolling} & {M.~hills} & {Hills} & {Mtn} \\
  & & & & &
  & {$(168)$} & {$(26)$} & {$(32)$} & {$(50)$} & {$(14)$} \\
\midrule
\rowcolor{black!5}
Oracle naive ($v_\text{true}$)
  & {--} & {--} & {--} & {--} & {--}
  & {--} & {--} & 3.4 & 5.1 & 11.5 \\
\rowcolor{black!5}
Oracle corrections
  & {\cmark} & {tuned} & {4-dur} & {cond} & {--}
  & {--} & {--} & 3.3 & 4.7 & 12.6 \\
\midrule
Moving average
  & {--} & {--} & {--} & {--} & {--}
  & {--} & 7.2 & 10.0 & 14.8 & 33.6 \\
Constant-power prior
  & {\cmark} & {--} & {--} & {--} & {--}
  & {--} & 9.7 & 11.6 & 16.1 & 34.8 \\
Behavioural prior
  & {\cmark} & {param} & {--} & {--} & {--}
  & {--} & 8.8 & 10.3 & 12.4 & 21.6 \\
\quad + single correction
  & {\cmark} & {param} & {global} & {--} & {--}
  & {--} & 7.4 & 9.8 & 13.1 & 18.8 \\
\quad + split correction
  & {\cmark} & {param} & {split} & {--} & {--}
  & {--} & 7.3 & 9.9 & 12.9 & 17.3 \\
\quad + profile-conditional
  & {\cmark} & {param} & {split} & {cond} & {--}
  & {--} & 7.3 & 9.9 & 12.9 & 17.8 \\
\quad + adaptive $\alpha$/$\kappa$
  & {\cmark} & {tuned} & {split} & {cond} & {--}
  & {--} & 7.3 & 9.8 & 12.5 & 16.1 \\
\quad + 4-dur corrections
  & {\cmark} & {param} & {4-dur} & {cond} & {--}
  & {--} & 7.2 & 9.2 & 12.3 & 16.1 \\
\quad + both (4-dur + adaptive)
  & {\cmark} & {tuned} & {4-dur} & {cond} & {--}
  & {--} & 7.2 & 9.2 & 12.0 & 14.0 \\
\quad + time lerp ($f^2$)
  & {\cmark} & {tuned} & {4-dur} & {cond} & {--}
  & {--} & 7.0 & 9.1 & 12.0 & 15.1 \\
\bottomrule
\end{tabular}
\end{table}

The two oracle rows frame the discussion. The oracle naive (true
whole-ride average speed, no physics) shows the \emph{within-ride pace
variability floor}: the best any constant-speed model can achieve.
Hills are low (5.6) because speed oscillations between bergs and flats
cancel over mixed terrain; mountains are high (11.3) because the
sustained speed difference between climbing and descending does not
cancel. The oracle corrections (physics model with perfect correction
values from row~0) show what the full apparatus achieves when
convergence is instantaneous: 5.3 on hills (below the naive floor,
proving physics helps when calibrated) and 8.2 on mountains. The gap
between these oracles and the real estimators below is pure
convergence overhead --- the cost of learning the corrections online.

Several patterns stand out. The moving average is a strong baseline on
flat rides (MAE 6.1\,min on pancake) but collapses on mountains
(29.9\,\si{\minute}) because it cannot anticipate upcoming terrain.
The constant-power physics prior is actually \emph{worse} than the
moving average everywhere (31.4 vs 29.9 on mountains, 7.3 vs 6.1 on
pancake): its descent predictions are so wrong that the extra gradient
information hurts. Adding the behavioral adjustments ($\alpha$,
$\kappa$) fixes the descent shape and recovers 5.7\,\si{\minute} on
mountains (25.7), but without correction the model still has the right
\emph{shape} and the wrong \emph{level} (MPE~$= -15$\,\%).

A single global integral correction cuts the mountain error nearly in
half (15.9), and splitting it into separate climb and descent channels
gives a further 1.6\,\si{\minute} (14.3) by letting each terrain class
converge independently.

The profile-conditional relevance filter is the key architectural
contribution. On mountain routes (identified at ride start from the
PCS climb-score classification) it gates the integral corrections to
learn only from segments whose predicted traversal time exceeds
120\,\si{\second}, reducing mountain MAE by 19\,\% (11.6) without
touching any other bucket. On non-mountain routes it is a no-op,
preserving the split correction's full learning rate on terrain where
every short climb is the signal. This conditional design ensures
that the filter \emph{augments the baseline when the route profile
warrants it, and falls back to the uncorrected split when it does
not} --- matching the ``never worse'' principle of
Section~\ref{sec:baseline}.

The power-curve rows show the additional benefit when a power meter is
available. A per-ride empirical $\pi(g)$ curve, trust-ramped against
the parametric fallback, improves hills (best at 14.5) and mountains
(10.6). A cached cross-ride profile improves mountains further (9.9)
at the cost of requiring historical data. In both cases the power
meter is optional: the no-power variant at 11.6\,\si{\minute} is
already a 61\,\% reduction from the moving-average baseline.


\subsection{Where in the ride does climbing appear?}
\label{sec:centroid}

The improvement over the moving average depends not only on how much
climbing a route contains, but on \emph{where} the climbing sits. We
define the \emph{climb centroid} as the climb-score-weighted centre of
mass of climbing position, expressed as a fraction of route distance:
\begin{equation}
\label{eq:centroid}
\gamma \;=\;
\frac{\sum_j s_j \cdot \bar{d}_j}{D \cdot \sum_j s_j},
\end{equation}
where $s_j$ is the PCS climb score of climb~$j$, $\bar{d}_j$ is its
midpoint distance, and $D$ is total route distance. A value of
$\gamma < 0.4$ means climbing is front-loaded; $0.4 \le \gamma \le
0.6$ means distributed; $\gamma > 0.6$ means back-loaded.

Table~\ref{tab:centroid} shows the MAE improvement of the
4-dur + adaptive estimator over the moving average on the full
737-ride corpus, partitioned by centroid bucket and difficulty.

\begin{table}[ht]
\centering
\caption{MAE improvement (minutes, higher = better) of the
4-dur + adaptive estimator over the moving average, by climb
centroid and route difficulty ($n = 737$ rides).}
\label{tab:centroid}
\sisetup{round-mode=places, round-precision=1}
\begin{tabular}{@{} l rrr rrr @{}}
\toprule
& \multicolumn{3}{c}{Hills ($n = 132$)}
& \multicolumn{3}{c}{Mountains ($n = 29$)} \\
\cmidrule(lr){2-4} \cmidrule(lr){5-7}
Centroid
  & {$n$} & {MA} & {$\Delta$}
  & {$n$} & {MA} & {$\Delta$} \\
\midrule
Front ($\gamma < 0.4$)
  & 24 & 22.8 & {$+5.9$}
  & 9  & 38.4 & {$+25.3$} \\
Mid ($0.4$--$0.6$)
  & 83 & 12.7 & {$+2.1$}
  & 19 & 31.2 & {$+16.8$} \\
Back ($\gamma > 0.6$)
  & 25 & 14.1 & {$+2.6$}
  & 1  & 34.9 & {$+19.4$} \\
\bottomrule
\end{tabular}
\end{table}

Two mechanisms interact. The physics model provides gradient
awareness from the first prediction, regardless of where the climbing
sits --- this is pure \emph{route knowledge} that the moving average
lacks. The integral corrections provide \emph{calibration} of that
knowledge, but they need observations to converge. The interplay
between these two mechanisms explains the centroid pattern.

\paragraph{Physics model quality.}
The final correction values $\Delta_c$ and $\Delta_d$ reveal how much
systematic bias the physics model carries before calibration. Across
the corpus, the median $\Delta_c$ at ride end is 0.98--1.02 on minor
hills (the model is nearly unbiased), 0.83--0.96 on hills (5--17\,\%
overestimate of climb speed), and 0.78--0.81 on mountains (a
$\sim$20\,\% overestimate). The model consistently overpredicts both
climb speed (rider pushes less hard than the behavioral model
assumes) and descent speed (rider brakes more than modelled). These
biases are what the integral corrections exist to absorb.

\paragraph{Correction convergence.}
On front-loaded mountain rides, the correction at the centroid
position is still 10.5\,\% away from its final value
($\Delta_c = 0.72$ vs 0.78). On mid-loaded mountain rides, the gap
shrinks to 1.5\,\% ($\Delta_c = 0.80$ vs 0.80). The front-loaded
case is ``physics alone'': the model is doing the heavy lifting with
a 20\,\% bias that the correction has not yet absorbed. On mountains
this still produces a $+$24\,\si{\minute} improvement because the
moving average's error is so catastrophic ($\sim$38\,\si{\minute})
that even an uncalibrated physics model is vastly better. On hills
the picture reverses: the physics model's 5--17\,\% bias is
comparable in magnitude to the moving average's error, so the
uncalibrated model can be slightly \emph{worse} ($-$1.8\,\si{\minute}
on front-loaded hills).

Mid-loaded routes are the sweet spot: enough distance has been
observed for the corrections to converge (gap $\sim$1--3\,\%), and
enough distance remains for the calibrated model to out-predict the
moving average. This is where the combination of physics model and
integral correction is strongest, producing $+$3.9\,\si{\minute}
on hills and $+$16.4\,\si{\minute} on mountains.

Back-loaded routes show the smallest improvement ($+$0.9 on hills,
$+$19.9 on mountains) simply because little distance remains when
the climbing begins --- the moving average's error window shrinks
faster than the model's advantage can accumulate. Both corrections
are fully converged, but there is not enough remaining ride for the
improvement to compound.


\subsection{Where does the error come from?}
\label{sec:oracle}

To understand the limits of any gradient-aware estimator, we
construct three oracle baselines that cheat with information
unavailable in practice:

\begin{itemize}
\item \textbf{Oracle naive.} Predict remaining time as
  $\text{remaining distance} / v_{\text{true}}$, where
  $v_{\text{true}}$ is the ride's actual average moving speed. No
  physics, no corrections, no gradient awareness --- just the perfect
  constant-speed projection. This measures the \emph{within-ride pace
  variability floor}: how much the rider's speed deviates from their
  own average throughout the ride.
\item \textbf{Oracle corrections.} Use the physics model with
  duration-bucketed corrections, but pre-set all four correction
  factors to their final end-of-ride values from row~0. No
  convergence delay, perfect calibration from the start.
\item \textbf{Naive with prior.} The constant-speed projection at the
  rider's configured $v_\text{flat}$ (\SI{29}{\km\per\hour}). This is
  what the rider gets if we strip away all gradient awareness but keep
  the prior.
\end{itemize}

\begin{table}[ht]
\centering
\caption{Oracle and ceiling analysis: MAE (minutes) by route
difficulty. The gap between oracle corrections and the real estimator
is the convergence overhead --- the cost of learning the corrections
online rather than knowing them in advance.}
\label{tab:oracle}
\sisetup{round-mode=places, round-precision=1}
\begin{tabular}{@{} l S[table-format=2.1] S[table-format=2.1] S[table-format=2.1] S[table-format=2.1] S[table-format=2.1] @{}}
\toprule
& {M.~hills} & {Hills} & {Mtn} \\
& {$(89)$} & {$(132)$} & {$(29)$} \\
\midrule
Oracle naive ($v_\text{true}$, no physics)
  & 3.4 & 5.1 & 11.5 \\
Oracle corrections (physics, perfect $\Delta$)
  & 3.3 & 4.7 & 12.6 \\
\midrule
4-dur + adaptive (real estimator)
  & 9.2 & 12.0 & 14.0 \\
Moving average
  & 10.0 & 14.8 & 33.6 \\
\bottomrule
\end{tabular}
\end{table}

Three findings emerge from Table~\ref{tab:oracle}.

\paragraph{The within-ride variability floor.}
Even knowing the rider's true average speed perfectly, the oracle
naive achieves only 5.1\,\si{\minute} on hills and 11.5 on
mountains --- not zero. The residual is from pace variation
\emph{within} the ride: fatigue in the second half, headwind on the
return leg, group dynamics, variable effort across different bergs.
No constant-speed model can go below this floor.

\paragraph{Physics helps on hills, hurts on mountains.}
On hills, oracle corrections (4.7) beat oracle naive (5.1) by
0.4\,\si{\minute} --- per-gradient predictions outperform the
aggregate when perfectly calibrated. On mountains, however, oracle
corrections (12.6) are \emph{worse} than oracle naive (11.5). This
reversal, absent in the smaller 290-ride corpus, appears because the
expanded mountain set (29 rides) includes more diverse terrain where
the ratio table's per-gradient shape errors are correlated and do not
cancel. The naive model's error-cancellation property is more
valuable than per-segment accuracy on these rides. This finding
suggests that the behavioral parameters ($\alpha = 0.25$,
$\kappa = 0.45$) tuned on the original 14 mountain rides do not
generalise to the broader mountain corpus and need retuning --- or
that the mountain category itself is too heterogeneous for a single
parameter set.

\paragraph{Convergence remains the dominant cost on hills.}
The gap between oracle corrections and the real estimator is
7.3\,\si{\minute} on hills (4.7 vs 12.0). On mountains the gap is
only 1.4\,\si{\minute} (12.6 vs 14.0), but this is misleading:
the oracle corrections themselves are already worse than the naive
floor, so the physics model's ceiling on mountains is set by its
shape quality, not by convergence speed. The practical implication
differs by terrain: on hills, faster convergence is the lever; on
mountains, better shape (retuned $\alpha$/$\kappa$ or the empirical
power curve) is the lever.


\subsection{When does physics help?}
\label{sec:when-helps}

Table~\ref{tab:convergence} decomposes the MAE by ride fraction,
revealing when the physics model earns its keep.

\begin{table}[ht]
\centering
\caption{MAE (minutes) computed from different starting fractions.
Later starts exclude the convergence period and show steady-state
prediction quality.}
\label{tab:convergence}
\sisetup{round-mode=places, round-precision=1}
\begin{tabular}{@{} l l S[table-format=2.1] S[table-format=2.1] S[table-format=2.1] @{}}
\toprule
& {Start from}  & {Pancake} & {Hills} & {Mtn} \\
\midrule
\multirow{3}{*}{Moving average}
& {2\,\%}  & 7.3 & 17.7 & 31.7 \\
& {25\,\%} & 4.8 & 12.2 & 26.0 \\
& {50\,\%} & 3.1 &  9.1 & 20.1 \\
\midrule
\multirow{3}{*}{Combined}
& {2\,\%}  & 7.1 & 15.1 & 12.1 \\
& {25\,\%} & 5.7 & 11.0 & 11.1 \\
& {50\,\%} & 3.8 &  8.4 & 11.4 \\
\bottomrule
\end{tabular}
\end{table}

The physics model solves exactly one problem: predicting speed on
terrain the rider has not yet reached. That problem is catastrophic
on mountains (the moving average is 26\,\si{\minute} off even at
25\,\% of the ride; the physics model is 11), modest on hills in the
first quarter (15.1 vs 17.7), and irrelevant on flat rides and
late in any ride where the moving average has converged.

On hills, after 25\,\% of the ride ($\sim$30--45\,\si{\minute} of
typical riding), the moving average and the physics model converge to
similar performance (12.2 vs 11.0). The corrections have settled, and
the remaining error is dominated by within-ride pace variability ---
group dynamics, effort choices on individual bergs, wind shifts ---
that no gradient-aware model can predict. The physics model's value
on hills is concentrated in the first quarter, where it provides a
2--3\,\si{\minute} head start over the moving average.

On mountains, the physics model provides a
$\sim$15\,\si{\minute} advantage at every stage of the ride, because
the sustained speed difference between climbing and descending never
cancels in the way that short berg-and-flat oscillations do. A rider
50\,\% into a mountain ride sees a moving-average ETA that is
20\,\si{\minute} wrong; the physics model is 11. This is the
scenario where gradient awareness is not an optimisation but a
necessity.

The ``never worse'' property holds throughout: on pancake rides the
physics model is within 1\,\si{\minute} of the moving average at
every fraction, confirming that the architecture degrades gracefully
when there is nothing to predict.

For a device-side implementation, the practical message is: the
physics model is essential equipment for mountain rides and useful
early insurance on hilly rides. On flat rides it is harmless ballast.
A rider who only rides flat terrain gains little from gradient-aware
ETA; a rider who occasionally ventures into the mountains gains a
great deal. The route-conditional architecture ensures both riders
get the best available prediction without configuration.


\section{Summary}

Method~D is a deliberately simple ETA estimator built around two
ideas. First, a physics-based speed-ratio table $r(g) = v(g) /
v_{\text{flat}}$ derived from a Cardano-solved force balance, with
small behavioral adjustments that make the curve shape-stable across
different rider operating points. Second, a cumulative integral
correction over observed climb and descent rows that absorbs any
steady bias the physics model cannot predict. A brief warmup period
uses a prior-free $v_{\text{flat}}$ observer to recalibrate the
ratio table at the rider's actual operating point, removing most of
the sensitivity to an initial $v_{\text{flat}}$ guess. The
time-to-go computation reduces to a dot product between the rider's
inverse-ratio vector and a histogram of remaining route length
binned by gradient, which is cheap to evaluate, trivial to update
incrementally, and small enough to fit comfortably in a bike
computer's memory footprint.

The method does not attempt to model fatigue, wind, group dynamics,
or mid-ride state changes. Instead it relies on the cumulative
integrals to absorb whatever steady bias these unmodeled effects
introduce, and on the ride being long enough for the expanding
window to dominate minor within-ride variations. For the rider's own
repeated rides this assumption holds comfortably; for rides with
violent mid-ride state changes (bonking, severe weather shifts) the
method's accuracy degrades gracefully rather than catastrophically.

\appendix

\section{Full metrics by route difficulty}
\label{app:metrics}

Table~\ref{tab:full-flat}--\ref{tab:full-mtn} report the complete
metric set for each estimator variant, partitioned by PCS-inspired
route difficulty. MAE and RMSE are in minutes; MPE and MAPE are
percentages of actual remaining time; Settle is the moving time in
minutes until the ETA prediction stabilises within $\pm$5\,\% of the
final value.

\begin{table}[ht]
\centering
\caption{Full metrics --- flat rides ($n = 226$).}
\label{tab:full-flat}
\sisetup{round-mode=places, round-precision=1}
\begin{tabular}{@{} l S[table-format=1.1] S[table-format=1.1] S[table-format=2.1] S[table-format=2.1] S[table-format=2.1] @{}}
\toprule
Estimator & {MAE} & {RMSE} & {MPE\,\%} & {MAPE\,\%} & {Settle} \\
\midrule
Moving average    & 7.1  & 8.7  & 0.1   & 9.2  & 25.4 \\
Constant-power    & 8.7  & 6.7  & -8.8  & 11.5 & 60.1 \\
Behavioural prior & 7.8  & 5.9  & -5.0  & 9.8  & 52.4 \\
+ single corr.    & 7.3  & 8.3  & -1.3  & 9.7  & 26.1 \\
+ split corr.     & 6.6  & 7.5  & -1.8  & 8.7  & 21.9 \\
+ relevance       & 6.7  & 7.6  & -0.5  & 8.6  & 24.9 \\
+ ride power       & 6.7  & 7.6  & -0.6  & 8.5  & 25.0 \\
\cmidrule(lr){1-6}
+ cached profile  & 6.7  & 7.6  & -0.7  & 8.4  & 24.0 \\
\bottomrule
\end{tabular}
\end{table}

\begin{table}[ht]
\centering
\caption{Full metrics --- hilly rides ($n = 50$).}
\label{tab:full-hills}
\sisetup{round-mode=places, round-precision=1}
\begin{tabular}{@{} l S[table-format=2.1] S[table-format=2.1] S[table-format=2.1] S[table-format=1.1] S[table-format=3.1] @{}}
\toprule
Estimator & {MAE} & {RMSE} & {MPE\,\%} & {MAPE\,\%} & {Settle} \\
\midrule
Moving average    & 17.5 & 15.0 & 1.5   & 8.8  & 54.9  \\
Constant-power    & 21.3 & 15.4 & 0.4   & 10.3 & 147.8 \\
Behavioural prior & 19.4 & 12.4 & -6.4  & 8.9  & 138.8 \\
+ single corr.    & 14.9 & 13.1 & -1.2  & 7.7  & 55.2  \\
+ split corr.     & 14.9 & 12.3 & -0.9  & 7.8  & 64.2  \\
+ relevance       & 15.1 & 11.9 & 2.0   & 8.1  & 56.7  \\
+ ride power       & 14.5 & 11.8 & 1.9   & 7.2  & 53.4  \\
\cmidrule(lr){1-6}
+ cached profile  & 14.7 & 12.1 & 1.5   & 7.4  & 57.5  \\
\bottomrule
\end{tabular}
\end{table}

\begin{table}[ht]
\centering
\caption{Full metrics --- mountain rides ($n = 14$).}
\label{tab:full-mtn}
\sisetup{round-mode=places, round-precision=1}
\begin{tabular}{@{} l S[table-format=2.1] S[table-format=2.1] S[table-format=3.1] S[table-format=2.1] S[table-format=3.1] @{}}
\toprule
Estimator & {MAE} & {RMSE} & {MPE\,\%} & {MAPE\,\%} & {Settle} \\
\midrule
Moving average    & 29.9 & 36.0 & 2.2    & 16.0 & 33.9  \\
Constant-power    & 31.4 & 31.2 & 14.6   & 17.3 & 196.4 \\
Behavioural prior & 25.7 & 31.5 & -15.2  & 15.2 & 211.2 \\
+ single corr.    & 15.9 & 18.4 & 2.6    & 8.7  & 23.3  \\
+ split corr.     & 14.3 & 18.9 & 0.3    & 7.7  & 14.5  \\
+ relevance       & 11.6 & 15.5 & -0.6   & 6.1  & 26.7  \\
+ ride power       & 10.6 & 13.9 & -4.3   & 7.8  & 26.4  \\
\cmidrule(lr){1-6}
+ cached profile  & 9.9  & 14.1 & -3.2   & 6.1  & 22.0  \\
\bottomrule
\end{tabular}
\end{table}

\end{document}
